% ============================================================
% sec/05analisis_tecnico.tex — TRACK A
%
% DEBE CONTENER, sin excepción:
%   - Análisis exploratorio PROFUNDO
%   - Diccionario de datos
%   - Calidad y estrategias de imputación JUSTIFICADAS
%   - Detección de valores atípicos
%   - Desbalance de clases
%   - Prevención EXPLÍCITA de data leakage
%
% El data leakage confirmado cuesta -15 puntos y corrección
% obligatoria. Es el riesgo que puede invalidar todo el trabajo.
% ============================================================
\section{Análisis Técnico del Track}
\label{sec:analisistecnico}

\pc{Una oración nombrando el track aprobado y la fecha de aprobación.}

\subsection{Análisis exploratorio}
\label{subsec:eda}
\pc{Número de registros, número de características, distribución del
target, cobertura temporal, y la estructura de agrupamiento de los
datos. Distinguir lo MEDIDO de lo PROPUESTO: nada propuesto debe
leerse como resultado.}

\subsection{Diccionario de datos}
\label{subsec:diccionario}
\begin{table}[!t]
\centering
\caption{Diccionario de datos (extracto).}
\label{tab:diccionario}
\footnotesize
\begin{tabularx}{\columnwidth}{@{}P{1.7cm} P{1.2cm} P{1.6cm} Y@{}}
\toprule
\textbf{Variable} & \textbf{Tipo} & \textbf{Origen} &
\textbf{Significado / rango / riesgo}\\
\midrule
\pc{var} & \pc{tipo} & \pc{origen} & \pc{descripción}\\
\bottomrule
\end{tabularx}
\end{table}

\pc{Si el diccionario completo no cabe, incluir un extracto
representativo y remitir al repositorio para la versión íntegra.}

\subsection{Calidad de los datos}
\label{subsec:calidad}
\pc{Qué se verificó y con qué resultado: duplicados, consistencia de
esquema, rangos imposibles, consistencia lógica entre campos,
cardinalidad de las categóricas.}

\subsection{Datos faltantes e imputación}
\label{subsec:imputacion}

\pc{Distinguir los mecanismos de ausencia. Tratarlos igual sería un
error, y la estrategia debe justificarse con literatura, no con
costumbre.}

\noindent\textbf{Ausencia estructural.} \pc{Campos ausentes porque el
dominio los define así. No son datos faltantes en sentido estadístico
y no se imputan: se codifican semánticamente. Justificar con
referencia.}

\noindent\textbf{Ausencia inesperada.} \pc{Campos requeridos ausentes
por corrupción o error de procesamiento. Procedimiento: medir la
frecuencia, decidir según su magnitud, y documentar el criterio.}

\noindent\textbf{Estrategia adoptada.} \pc{Qué se hace con cada
mecanismo y por qué. Ninguna variable se imputa automáticamente sin
estudiar antes el mecanismo de ausencia.}

\subsection{Valores atípicos}
\label{subsec:atipicos}
\pc{Método de detección, umbral, y qué se hace con lo detectado. Un
valor extremo puede ser un atípico VÁLIDO y no un error: solo se
elimina lo demostradamente corrupto, y la regla de exclusión se declara
antes de mirar el conjunto de prueba.}

\subsection{Desbalance de clases}
\label{subsec:desbalance}
\pc{Proporción medida sobre la representación final, no sobre una
aproximación. Estrategia elegida y su justificación con literatura.
Si la salida del modelo es una probabilidad, discutir el efecto de la
corrección sobre la calibración antes de aplicarla.}

\subsection{Prevención explícita de data leakage}
\label{subsec:leakage}

\pc{Este es el riesgo metodológico propio del Track A y tiene un
criterio dedicado en la rúbrica.}

\subsubsection{Lista negra de información prohibida}
\pc{Enumerar todo campo que no puede entrar al vector de
características, y por qué en el momento real de la decisión esa
información no existiría todavía.}

\subsubsection{Taxonomía de riesgo}
\begin{table}[!t]
\centering
\caption{Vectores de fuga: riesgo, evidencia y mitigación.}
\label{tab:leakage}
\footnotesize
\begin{tabularx}{\columnwidth}{@{}P{1.9cm} Y P{2.2cm}@{}}
\toprule
\textbf{Tipo} & \textbf{Riesgo y evidencia medida} &
\textbf{Mitigación}\\
\midrule
\pc{tipo} & \pc{riesgo + evidencia} & \pc{mitigación}\\
\pc{tipo} & \pc{riesgo + evidencia} & \pc{mitigación}\\
\pc{tipo} & \pc{riesgo + evidencia} & \pc{mitigación}\\
\bottomrule
\end{tabularx}
\end{table}

\subsubsection{Verificación automática}
\pc{Qué prueba automática falla si un campo prohibido entra al modelo
o si una unidad independiente aparece en dos particiones. Una política
sin prueba que la verifique es una intención, no una mitigación.}